\preliminar{Glosario}

\begin{description}[style=nextline,leftmargin=1.27cm,labelindent=0pt,%
  labelsep=0.4em,font=\bfseries,itemsep=0.25\baselineskip]
  \item[Cotización] Propuesta económica que detalla el precio, los
        servicios incluidos y las condiciones de un paquete turístico.
  \item[Cupo] Cantidad de lugares disponibles que quedan en una salida
        turística luego de descontar las reservas confirmadas.
  \item[Disponibilidad] Condición que indica si un paquete o una salida
        puede venderse en una fecha determinada.
  \item[GDS] Sistema global de distribución que permite consultar
        horarios, tarifas y disponibilidad de varios proveedores
        turísticos desde un mismo punto de acceso.
  \item[Itinerario] Secuencia de actividades, traslados y servicios
        previstos durante un viaje.
  \item[Paquete turístico] Combinación de transporte, alojamiento y
        actividades que se ofrece al cliente a un precio determinado.
  \item[Reserva] Registro de un cliente y de los pasajeros que ocuparán
        los cupos de una salida turística.
  \item[Sistema web] Aplicación a la que se accede mediante un navegador
        y que centraliza la información en un servidor.
  \item[Trazabilidad] Capacidad de seguir el historial y el estado de
        cada reserva o cambio realizado en el sistema.
\end{description}

\section*{Acrónimos}

\begin{description}[style=nextline,leftmargin=1.27cm,labelindent=0pt,%
  labelsep=0.4em,font=\bfseries,itemsep=0.15\baselineskip]
  \item[GDS] \emph{Global Distribution System}
  \item[INE] Instituto Nacional de Estadística
  \item[ISO] \emph{International Organization for Standardization}
  \item[OWASP] \emph{Open Worldwide Application Security Project}
  \item[TPS] \emph{Transaction Processing System}
  \item[UPDS] Universidad Privada Domingo Savio
\end{description}
